
\subsection{\texorpdfstring{\Spherical}{Spherical} Isometry}\label{sec:sfin}

The \texttt{SFIN} transformation 
was derived for finite element analysis %from the investigation
%\FCF{appears in the study}
by \cite{crisfield1999objectivity}.
% This procedure is invariant to node-numbering. 
For two-node elements, this transformation effects a 
spherical linear interpolation of the rotation field
which furnishes
%\FCF{It turns out that} a spherical linear interpolation 
 a geodesic on \(\mathrm{SO}(3)\) 
and provides the most natural generalization of 
a %classical
linear interpolation for the rotation manifold \citep{shoemake1985animating}.
%For two rotations \(\FrameRotation_I\) and \(\FrameRotation_J\), the interpolating rotation \(\FrameRotation^h\left(\xi\right)\)  is:
% $$
% \FrameRotation^h\left(\xi\right)% ; \FrameRotation_I, \FrameRotation_J\right)
% =\FrameRotation_I \exp \left(\frac{\xi}{L} \log \left(\FrameRotation_I^\mathrm{t} \FrameRotation_J\right)\right)
% $$
%
\cite{jelenić1999geometrically} use the concept 
for a finite element formulation in the context of dynamics.
The study used a Petrov-Galerkin methodology that resulted
in a non-symmetric stiffness,
 and does not distinguish between element-independent operations and element-dependent operations.
% \FCF{without distinction of element-independent from element-dependent operations}.
The external application of the transformation in the current study to a wrapped element
% the elements
%like that
%of \cite{ibrahimbegović1995computational} and \cite{simo1986threedimensional}
imparts %to them 
the benefits described in \cite{crisfield1999objectivity},
while enforcement of the Bubnov-Galerkin condition, % as proposed here,
ensures symmetry of the tangent stiffness at equilibrium.

The procedure for forming \(\bm{R}\) begins
with the selection of two node indices \(I\) and \(J\) for an \(n\)-noded element:
%
\begin{equation*} 
I=\operatorname{floor}\left(\frac{1}{2}(n+1)\right)  
\quad \text { and } \quad  
J=\operatorname{floor}\left(\frac{1}{2}(n+2)\right).
\end{equation*}
%
The rotations \(\FrameRotation_I\) and \(\FrameRotation_J\)
at nodes $I$ and $J$
%, respectively,
are used to form an intermediate vector $\bm{\vartheta}$ with the formula:
\[\bm{\vartheta} = \operatorname{Log}\FrameRotation_I^{\mathrm{t}} \FrameRotation_J.\]
With this vector, the coordinate rotation is given by:
%
\begin{equation*} 
\left.\begin{array}{rl}
\bm{R} &= \FrameRotation_I \operatorname{Exp} \left(\frac{1}{2} \bm{\vartheta}\right). \\
\end{array}\right.
\end{equation*}
%
The nodal spin matrices \(\bm{W}_a\) %each
comprise two $3\times 3$ blocks each with the form:
%
\begin{equation*} 
\bm{W}_a = \left[\bm{0} ~~ \bm{W}_{a\scriptscriptstyle{\Lambda}}\right]
\end{equation*}
%
where \(\bm{W}_{a\scriptscriptstyle{\Lambda}} = \bm{0}\) for \(a \notin \{I,J\}\) and otherwise:
%
\begin{equation*} 
\begin{aligned}
\bm{W}_{I \Lambda} &= \frac{1}{2}\left(\bm{1} + \frac{1}{\vartheta} \tan \frac{\vartheta}{4} \, \bm{\vartheta}^{\times}\right) \\
\end{aligned}
\qquad\text{ and }\qquad
\begin{aligned}
\bm{W}_{J \Lambda} &=\frac{1}{2}\left(\bm{1} -\frac{1}{\vartheta} \tan \frac{\vartheta}{4} \, \bm{\vartheta}^{\times}\right)
\end{aligned}
\end{equation*}
with \(\vartheta = \|\bm{\vartheta}\|\).
% \begin{equation*} 
% \operatorname{Sfin}\left(\FrameRotation_I, \FrameRotation_J, \xi\right)=\FrameRotation_I \exp \left(\frac{\xi}{L} \log \left(\FrameRotation_I^\mathrm{t} \FrameRotation_J\right)\right)
% \end{equation*}

\begin{ambox}[\texttt{Sfin} Isometry]{box:sfin}
    
\end{ambox}
